% Options for packages loaded elsewhere
\PassOptionsToPackage{unicode}{hyperref}
\PassOptionsToPackage{hyphens}{url}
\PassOptionsToPackage{dvipsnames,svgnames,x11names}{xcolor}
\documentclass[
  11pt,
  a4paper,
]{article}
\usepackage{xcolor}
\usepackage[margin=1in]{geometry}
\usepackage{amsmath,amssymb}
\setcounter{secnumdepth}{-\maxdimen} % remove section numbering
\usepackage{iftex}
\ifPDFTeX
  \usepackage[T1]{fontenc}
  \usepackage[utf8]{inputenc}
  \usepackage{textcomp} % provide euro and other symbols
\else % if luatex or xetex
  \usepackage{unicode-math} % this also loads fontspec
  \defaultfontfeatures{Scale=MatchLowercase}
  \defaultfontfeatures[\rmfamily]{Ligatures=TeX,Scale=1}
\fi
\usepackage{lmodern}
\ifPDFTeX\else
  % xetex/luatex font selection
  \ifLuaTeX
    \usepackage{luaotfload}
    \directlua{luaotfload.add_fallback("mainfontfallback",{
      "NotoSerifBengali:mode=harf"
    })}
  \fi
  \setmainfont[,RawFeature={fallback=mainfontfallback}]{Latin Modern Roman}
  \ifLuaTeX
    \usepackage{luaotfload}
    \directlua{luaotfload.add_fallback("sansfontfallback",{
      "NotoSansBengali:mode=harf"
    })}
  \fi
  \setsansfont[,RawFeature={fallback=sansfontfallback}]{Latin Modern Sans}
  \ifLuaTeX
    \usepackage{luaotfload}
    \directlua{luaotfload.add_fallback("monofontfallback",{
      "NotoSansBengali:mode=harf"
    })}
  \fi
  \setmonofont[,RawFeature={fallback=monofontfallback}]{DejaVu Sans Mono}
\fi
% Use upquote if available, for straight quotes in verbatim environments
\IfFileExists{upquote.sty}{\usepackage{upquote}}{}
\IfFileExists{microtype.sty}{% use microtype if available
  \usepackage[]{microtype}
  \UseMicrotypeSet[protrusion]{basicmath} % disable protrusion for tt fonts
}{}
\makeatletter
\@ifundefined{KOMAClassName}{% if non-KOMA class
  \IfFileExists{parskip.sty}{%
    \usepackage{parskip}
  }{% else
    \setlength{\parindent}{0pt}
    \setlength{\parskip}{6pt plus 2pt minus 1pt}}
}{% if KOMA class
  \KOMAoptions{parskip=half}}
\makeatother
\usepackage{longtable,booktabs,array}
\usepackage{calc} % for calculating minipage widths
% Correct order of tables after \paragraph or \subparagraph
\usepackage{etoolbox}
\makeatletter
\patchcmd\longtable{\par}{\if@noskipsec\mbox{}\fi\par}{}{}
\makeatother
% Allow footnotes in longtable head/foot
\IfFileExists{footnotehyper.sty}{\usepackage{footnotehyper}}{\usepackage{footnote}}
\makesavenoteenv{longtable}
\setlength{\emergencystretch}{3em} % prevent overfull lines
\providecommand{\tightlist}{%
  \setlength{\itemsep}{0pt}\setlength{\parskip}{0pt}}
% paper_preamble.tex — shared LaTeX preamble for the PCMA/SI paper set.
%
% Compiled with lualatex (NOT pdflatex) because several papers contain
% Bengali script: the CLL paper realises one cognitive state into
% English, Bengali and Banglish, so the Bengali output is EVIDENCE and
% cannot be transliterated away without destroying the point being made.
%
% Font selection is NOT done here. pandoc >= 3 has native
% `mainfontfallback` / `monofontfallback` variables that emit the
% luaotfload fallback registration in the correct position relative to
% unicode-math; doing it by hand in this header lands after unicode-math
% and breaks font loading outright. See tools/build_papers.sh.

\usepackage{amsmath}
\usepackage{amssymb}
\usepackage{mathtools}   % \xleftrightarrow etc. used in cll.md §5
\usepackage{booktabs}

% ---------------------------------------------------------------------
% Unicode symbols appearing in PROSE rather than inside $...$.
%
% The papers were written as markdown and read as markdown, so they use
% real characters — "the ω-boost", "α-child", "ΔR coherence" — instead of
% LaTeX macros. Under unicode-math those codepoints are math-mode-only
% and abort the build ("\symbballowed only in math mode"). Mapping each
% to \ensuremath keeps the markdown readable on GitHub AND compilable,
% which is the whole point of keeping one source.
% ---------------------------------------------------------------------
\usepackage{newunicodechar}

% Greek, lower case
\newunicodechar{α}{\ensuremath{\alpha}}
\newunicodechar{β}{\ensuremath{\beta}}
\newunicodechar{γ}{\ensuremath{\gamma}}
\newunicodechar{δ}{\ensuremath{\delta}}
\newunicodechar{ε}{\ensuremath{\varepsilon}}
\newunicodechar{η}{\ensuremath{\eta}}
\newunicodechar{θ}{\ensuremath{\theta}}
\newunicodechar{κ}{\ensuremath{\kappa}}
\newunicodechar{λ}{\ensuremath{\lambda}}
\newunicodechar{μ}{\ensuremath{\mu}}
\newunicodechar{π}{\ensuremath{\pi}}
\newunicodechar{ρ}{\ensuremath{\rho}}
\newunicodechar{σ}{\ensuremath{\sigma}}
\newunicodechar{τ}{\ensuremath{\tau}}
\newunicodechar{φ}{\ensuremath{\varphi}}
\newunicodechar{χ}{\ensuremath{\chi}}
\newunicodechar{ψ}{\ensuremath{\psi}}
\newunicodechar{ω}{\ensuremath{\omega}}
\newunicodechar{ν}{\ensuremath{\nu}}
\newunicodechar{ξ}{\ensuremath{\xi}}
\newunicodechar{ζ}{\ensuremath{\zeta}}
\newunicodechar{ι}{\ensuremath{\iota}}

% Greek, upper case
\newunicodechar{Δ}{\ensuremath{\Delta}}
\newunicodechar{Φ}{\ensuremath{\Phi}}
\newunicodechar{Σ}{\ensuremath{\Sigma}}
\newunicodechar{Ψ}{\ensuremath{\Psi}}
\newunicodechar{Ω}{\ensuremath{\Omega}}
\newunicodechar{Γ}{\ensuremath{\Gamma}}
\newunicodechar{Λ}{\ensuremath{\Lambda}}
\newunicodechar{Θ}{\ensuremath{\Theta}}
\newunicodechar{Π}{\ensuremath{\Pi}}

% Operators and relations
\newunicodechar{×}{\ensuremath{\times}}
\newunicodechar{·}{\ensuremath{\cdot}}
\newunicodechar{±}{\ensuremath{\pm}}
\newunicodechar{−}{\ensuremath{-}}
\newunicodechar{≈}{\ensuremath{\approx}}
\newunicodechar{≤}{\ensuremath{\leq}}
\newunicodechar{≥}{\ensuremath{\geq}}
\newunicodechar{≠}{\ensuremath{\neq}}
\newunicodechar{∈}{\ensuremath{\in}}
\newunicodechar{∂}{\ensuremath{\partial}}
\newunicodechar{∇}{\ensuremath{\nabla}}
\newunicodechar{√}{\ensuremath{\sqrt{}}}
\newunicodechar{∞}{\ensuremath{\infty}}
\newunicodechar{∝}{\ensuremath{\propto}}
\newunicodechar{∑}{\ensuremath{\sum}}
\newunicodechar{∫}{\ensuremath{\int}}
\newunicodechar{ℝ}{\ensuremath{\mathbb{R}}}
\newunicodechar{ℕ}{\ensuremath{\mathbb{N}}}
\newunicodechar{‖}{\ensuremath{\|}}
\newunicodechar{⊕}{\ensuremath{\oplus}}
\newunicodechar{⊗}{\ensuremath{\otimes}}
\newunicodechar{∀}{\ensuremath{\forall}}
\newunicodechar{∃}{\ensuremath{\exists}}
\newunicodechar{¬}{\ensuremath{\neg}}
\newunicodechar{∧}{\ensuremath{\wedge}}
\newunicodechar{∨}{\ensuremath{\vee}}

% Arrows
\newunicodechar{→}{\ensuremath{\rightarrow}}
\newunicodechar{←}{\ensuremath{\leftarrow}}
\newunicodechar{↔}{\ensuremath{\leftrightarrow}}
\newunicodechar{↑}{\ensuremath{\uparrow}}
\newunicodechar{↓}{\ensuremath{\downarrow}}
\newunicodechar{⇒}{\ensuremath{\Rightarrow}}
\newunicodechar{⇐}{\ensuremath{\Leftarrow}}
\newunicodechar{⇔}{\ensuremath{\Leftrightarrow}}

% Super/subscripts written as literal characters
\newunicodechar{²}{\ensuremath{^{2}}}
\newunicodechar{³}{\ensuremath{^{3}}}
\newunicodechar{⁴}{\ensuremath{^{4}}}
\newunicodechar{⁵}{\ensuremath{^{5}}}
\newunicodechar{⁶}{\ensuremath{^{6}}}
\newunicodechar{₀}{\ensuremath{_{0}}}
\newunicodechar{₁}{\ensuremath{_{1}}}
\newunicodechar{₂}{\ensuremath{_{2}}}
\newunicodechar{₃}{\ensuremath{_{3}}}
\newunicodechar{ᵀ}{\ensuremath{^{\mathsf{T}}}}

% Status glyphs. The papers use emoji as status markers in tables
% (implemented / failed / warning). Mapped to typographic equivalents:
% a colour emoji font in lualatex is unreliable, and the MEANING is what
% carries — a reader needs to see "this one failed", not a red cross.
\newunicodechar{✅}{\ensuremath{\checkmark}}
\newunicodechar{✓}{\ensuremath{\checkmark}}
\newunicodechar{❌}{\ensuremath{\times}}
\newunicodechar{✗}{\ensuremath{\times}}
\newunicodechar{⚠}{\textbf{!}}
\newunicodechar{🧩}{\textbf{[piece]}}
\newunicodechar{📋}{\textbf{[list]}}
\newunicodechar{🟡}{\textbf{[partial]}}
\newunicodechar{🔴}{\textbf{[dead]}}
\newunicodechar{🟢}{\textbf{[live]}}
\newunicodechar{▼}{\ensuremath{\blacktriangledown}}
\newunicodechar{►}{\ensuremath{\blacktriangleright}}
\newunicodechar{□}{\ensuremath{\square}}
\newunicodechar{█}{\rule{0.6em}{0.9ex}}
\newunicodechar{•}{\textbullet}
\newunicodechar{†}{\textdagger}
\newunicodechar{‑}{-}
\newunicodechar{‐}{-}
\newunicodechar{…}{\ldots}
\newunicodechar{।}{\textbar}   % danda: rendered by the Bengali fallback in text
\newunicodechar{⁻}{\ensuremath{^{-}}}

% Zero-width formatting codepoints. These are invisible in an editor but
% are real characters to LaTeX: an emoji written as "⚠️" is TWO
% codepoints (U+26A0 U+FE0F), so declaring only the visible one leaves a
% variation selector with no definition and the build dies on a symbol
% nobody can see. Both are mapped to nothing.
\newunicodechar{^^^^fe0f}{}    % VARIATION SELECTOR-16
\newunicodechar{^^^^200d}{}    % ZERO WIDTH JOINER
\newunicodechar{^^^^200b}{}    % ZERO WIDTH SPACE

\usepackage{longtable}
\usepackage{array}

% Comparison matrices and cost ledgers are wide. Let them shrink to the
% text block rather than overflow the margin.
\usepackage{adjustbox}

% Code blocks: wrap long lines instead of running off the page. Several
% papers quote real C++ with long identifiers.
\usepackage{fvextra}
\fvset{breaklines=true,breakanywhere=true,fontsize=\small}

\setlength{\parskip}{0.5em}
\setlength{\parindent}{0pt}

\usepackage{fancyhdr}
\pagestyle{fancy}
\fancyhf{}
\fancyfoot[C]{\thepage}
\fancyhead[L]{\small\itshape Synthetic Intelligence — research lineage}
\fancyhead[R]{\small\itshape irfan@furylogic.com}
\renewcommand{\headrulewidth}{0.2pt}
\setlength{\headheight}{14pt}
\setlength{\headsep}{18pt}
\usepackage{bookmark}
\IfFileExists{xurl.sty}{\usepackage{xurl}}{} % add URL line breaks if available
\urlstyle{same}
\hypersetup{
  pdftitle={SRG: Synaptic Receptor Gating --- Dual-Port Cavities as an Alternative to Structural Fission},
  pdfauthor={Irfan Mahir --- Independent Researcher, Dhaka, Bangladesh --- irfan@furylogic.com},
  colorlinks=true,
  linkcolor={black},
  filecolor={Maroon},
  citecolor={Blue},
  urlcolor={[HTML]{0B5394}},
  pdfcreator={LaTeX via pandoc}}

\title{SRG: Synaptic Receptor Gating --- Dual-Port Cavities as an
Alternative to Structural Fission}
\author{Irfan Mahir --- Independent Researcher, Dhaka, Bangladesh ---
irfan@furylogic.com}
\date{}

\begin{document}
\maketitle

{
\hypersetup{linkcolor=}
\setcounter{tocdepth}{3}
\tableofcontents
}
\section{SRG: Synaptic Receptor Gating --- Dual-Port Cavities as an
Alternative to Structural
Fission}\label{srg-synaptic-receptor-gating-dual-port-cavities-as-an-alternative-to-structural-fission}

\textbf{Irfan Mahir} \emph{Independent Researcher, Dhaka, Bangladesh}
irfan@furylogic.com

\begin{quote}
\emph{``A neuron does not split when it receives conflicting input. It
receives it through different receptors.''}
\end{quote}

\begin{center}\rule{0.5\linewidth}{0.5pt}\end{center}

\subsection{Abstract}\label{abstract}

When excitatory and inhibitory waves arrive at a single cavity in a
wave-propagation substrate, they sum. Summation of opposed signals is
destructive interference: the contradictory structure carried by the two
waves cancels, and the plasticity updates they would have produced
cancel with it. A concept holding genuine contradictions therefore
learns nothing, because its updates annihilate.

\textbf{Resonant Bifurcation} solves this by splitting the cavity into
two. This paper specifies the alternative: keep one cavity and
\textbf{split its input ports}, exactly as a biological neuron uses
distinct receptor types --- AMPA for excitation, GABA for inhibition ---
on a single cell body.

Each cavity is given an \textbf{E-port} admitting only positive-polarity
tunnels and an \textbf{I-port} admitting only negative-polarity ones.
The two channels are integrated separately and combined only at readout,
so signed cancellation cannot occur during accumulation. Graph topology
stays fixed, no dynamic allocation is required, and the energy delivered
by a query is not divided between child nodes.

The mechanism is implemented and opt-in, and §6 states honestly what
that costs and when the alternative is preferable.

\begin{center}\rule{0.5\linewidth}{0.5pt}\end{center}

\subsection{1. The Destructive Interference
Problem}\label{the-destructive-interference-problem}

Consider a cavity receiving two tunnels, one of positive polarity and
one of negative. In a scalar cavity, both contribute to the same
displacement variable:

\[
x \;\leftarrow\; x + \Delta_{+} + \Delta_{-}
\]

When \(\Delta_{+} \approx -\Delta_{-}\), the cavity is left at rest and
the field concludes nothing happened. But something \emph{did} happen:
two structurally significant events arrived. The information was not
processed and discarded --- it was destroyed at the point of summation.

The consequence for learning is worse than the consequence for readout.
Hebbian consolidation is driven by the product of endpoint velocities.
If the velocity is cancelled, the product is near zero, and the tunnels
that delivered the conflicting evidence are \textbf{not} strengthened. A
concept that holds contradictions is therefore the one concept that
cannot learn about them.

This is the phenomenon that motivated structural fission.

\begin{center}\rule{0.5\linewidth}{0.5pt}\end{center}

\subsection{2. Dual-Port Cavities}\label{dual-port-cavities}

Rather than one intake, each cavity is given two:

\begin{itemize}
\tightlist
\item
  \textbf{E-port} --- accepts waves from positive-polarity tunnels only.
\item
  \textbf{I-port} --- accepts waves from negative-polarity tunnels only.
\end{itemize}

The cavity therefore carries four state variables instead of two:
\((x_E, v_E)\) and \((x_I, v_I)\).

The biological reading is direct. A cortical neuron does not maintain a
single undifferentiated input. Excitatory and inhibitory afferents
terminate on different receptor populations with different reversal
potentials, and the cell integrates them through separate conductances
before they jointly determine whether it fires.

\begin{center}\rule{0.5\linewidth}{0.5pt}\end{center}

\subsection{3. Mechanics}\label{mechanics}

Waves entering the E-port and I-port do not sum into a shared
displacement field during accumulation. They are integrated in separate
channels, and the observable displacement is derived only at
end-of-tick:

\[
x \;=\; x_E - x_I, \qquad v \;=\; v_E - v_I
\]

Equivalently, the I-port contribution is phase-inverted before it
reaches the observable --- a \(180^\circ\) shift --- so an inhibitory
signal carries an explicit negative functional sign rather than silently
erasing an excitatory one.

\textbf{The property that matters is where the subtraction happens.}
Both formulations produce the same observable \(x\). They differ in what
is available to plasticity: because Hebbian updates are computed from
the \emph{port} velocities rather than the observable, an excitatory
update lands in \(v_E\) products and an inhibitory one in \(v_I\)
products, and they cannot cancel. Structure accumulates in both channels
simultaneously.

With gating disabled, \(x_I\) and \(v_I\) remain zero and the observable
equals the E-port exactly, so the baseline dynamics are recovered
bit-for-bit. The mechanism is a strict extension.

\begin{center}\rule{0.5\linewidth}{0.5pt}\end{center}

\subsection{4. Comparison with Structural
Fission}\label{comparison-with-structural-fission}

Both mechanisms address the same failure. They differ in cost.

\begin{longtable}[]{@{}
  >{\raggedright\arraybackslash}p{(\linewidth - 4\tabcolsep) * \real{0.3333}}
  >{\raggedright\arraybackslash}p{(\linewidth - 4\tabcolsep) * \real{0.3333}}
  >{\raggedright\arraybackslash}p{(\linewidth - 4\tabcolsep) * \real{0.3333}}@{}}
\toprule\noalign{}
\begin{minipage}[b]{\linewidth}\raggedright
Property
\end{minipage} & \begin{minipage}[b]{\linewidth}\raggedright
Fission (RB+G)
\end{minipage} & \begin{minipage}[b]{\linewidth}\raggedright
Receptor gating (SRG)
\end{minipage} \\
\midrule\noalign{}
\endhead
\bottomrule\noalign{}
\endlastfoot
Graph topology & mutates --- new vertices spawned & \textbf{static} \\
Allocation & dynamic; parent/alias ledger required & none \\
Query energy & divided among children & \textbf{delivered whole to one
cavity} \\
State per concept & grows & fixed, 2× scalar \\
Expressive gain & two independently-addressable concepts & one concept,
two insulated pathways \\
Audit trail & explicit fission event with provenance & none --- no
structural event occurs \\
\end{longtable}

\textbf{Static topology is the practical advantage.} Fission requires
allocating vertices at runtime and maintaining a ledger mapping a parent
label to its children, so that later queries resolve correctly. Gating
requires neither: the graph after a thousand ticks is the graph that was
loaded.

\textbf{Undivided energy is the functional advantage.} Under fission,
driving a concept excites a parent whose energy then distributes across
children, so the amplitude reaching any single child is reduced. This
was a live concern for recall gain. Under gating, the drive reaches one
cavity at full amplitude and the tension is processed through two
insulated pathways inside it.

\begin{center}\rule{0.5\linewidth}{0.5pt}\end{center}

\subsection{5. What Gating Does Not
Provide}\label{what-gating-does-not-provide}

The comparison above would be dishonest without this section.

\textbf{No new concepts.} Fission produces two \emph{addressable}
entities: \texttt{metformin\_alpha} can acquire its own edges,
participate in its own analogies, and be reported separately. Gating
produces one entity with better internal bookkeeping. If the
contradiction is genuinely two distinct senses of a word, gating does
not separate them --- it only stops them destroying each other.

\textbf{No structural audit trail.} A fission event is a recorded
structural mutation with provenance, and downstream analysis can ask
\emph{when} and \emph{why} a concept split. Gating changes nothing
structurally, so there is nothing to record and nothing to explain.

\textbf{Doubled state.} Every cavity pays for two ports whether or not
it ever receives conflicting input. Fission pays only for concepts that
actually contradict.

The mechanisms are therefore \textbf{complementary rather than
competing}: gating is the cheap always-on defence against cancellation,
fission is the expensive response to a contradiction severe enough to
warrant two concepts.

\begin{center}\rule{0.5\linewidth}{0.5pt}\end{center}

\subsection{6. Implementation Status}\label{implementation-status}

Dual-port integration is implemented in the current substrate as the
\texttt{x\_E} / \texttt{v\_E} and \texttt{x\_I} / \texttt{v\_I} channels
on each cavity, with the observable derived at end-of-tick per §3.

It is \textbf{opt-in and off by default.} With SRG disabled the
inhibitory channel stays at zero and the substrate reproduces baseline
behaviour exactly, which is what makes the extension safe to carry. It
is not enabled by default because the fission path is the mechanism for
which reproducible multi-seed measurements exist; enabling gating by
default would change the baseline against which those results were
obtained.

An honest statement of the evidence: the cancellation problem of §1 is
real and observed, the fix in §3 is sound by construction, and a
head-to-head measurement of gating against fission on the same ontology
\textbf{has not been run}. Until it is, the comparison table in §4 is an
argument from mechanism, not a report of a result.

\begin{center}\rule{0.5\linewidth}{0.5pt}\end{center}

\subsection{7. Conclusion}\label{conclusion}

Receptor gating is the conservative answer to destructive interference:
it changes how a cavity integrates its inputs rather than how many
cavities exist. It costs one extra pair of state variables per cavity,
requires no dynamic allocation, and preserves the amplitude of a driven
query.

It does not replace fission, because it cannot create a concept. It
removes the reason a concept needed to be created in the \emph{common}
case --- where the conflict is signed input to one coherent idea rather
than two ideas sharing a name.

\begin{center}\rule{0.5\linewidth}{0.5pt}\end{center}

\subsection{References}\label{references}

\begin{itemize}
\tightlist
\item
  Mahir, I. \emph{RB+G: Gated Recursive Bifurcation with Analogical
  Bridging.} --- the fission mechanism this paper offers an alternative
  to.
\item
  Mahir, I. \emph{From Artificial to Synthetic Intelligence.} §5.4c ---
  the substrate carrying both mechanisms.
\item
  Mahir, I. \emph{Recursive Integrity Field Architecture (RIFA).} ---
  the contradiction metabolism that detects the tension in the first
  place.
\item
  Implementation: \texttt{src/physics.hpp} (\texttt{x\_E},
  \texttt{v\_E}, \texttt{x\_I}, \texttt{v\_I}; end-of-tick observable
  derivation).
\end{itemize}

\end{document}
